\chapter{Results and Evaluation}

\section{Frozen reference output}
The frozen \code{reference_v1} produces 1,414 union flags across 81 of the
85 stored series. The flag CSV has 1,399 company-share flags and
15 PX1 index flags; TBIDX has none. There are 198 high-volume flags, 1,237 return flags,
and 21 combined rows: $198+1,237-21=1,414$. Local news matches 548 flags;
866 have no match in the stored ilboursa file. The refined review queue has
ten spike rows and six decay rows. The original three sourced positives
still yield one primary-window hit, GIF, and misses for LSTR and SERVI.

\section{Protocol and event register}
The Phase 2 evaluation protocol was committed before candidate detector code
and expanded-label hit tables. It froze a development split before
1 January 2025 and a held-out split on and after that date. The primary
event window is $[-30,+5]$ calendar days around the sourced public date;
the descriptive secondary window is $[-60,+5]$. A positive enters the
denominator only if its source, public date, inclusion flag, and market-data
coverage are valid. The expanded register added 18 sourced rows. After the
source audit, those additions are 13 negative controls, three exploratory
rows, and two positives. Seven positive rows have dated source URLs. The evaluation source
metric counts six because it requires the headline-inclusion flag, which
is false for unlisted TSI. Of those six, five have a trading observation
in the primary window: four development and one held-out.

The audit aligned UADH with the 25 March 2026 suspension stated in its BVMT
notice. Syphax's 19 December 2023 anchor is a public announcement; the
11 December court ruling is recorded only as context. ADWYA's source does
not establish a takeover reason, so its suspension is exploratory and
excluded from negative-control counts. PGH and LNDOR share a notice;
control flags are reported with both separate and collapsed accounting.
These corrections were governed by a source-alignment rule committed before
results were recalculated, and no method's hit status changed.

\section{Uncertainty and random-placement comparison}
For $k$ hits out of $n$ covered events, let $\hat p=k/n$,
$z=1.959963984540054$, and $D=1+z^2/n$. The implemented 95\% Wilson
centre and half-width are
\[
 C=\frac{\hat p+z^2/(2n)}{D},\qquad
 H=\frac{z\sqrt{\hat p(1-\hat p)/n+z^2/(4n^2)}}{D}.
\]
The reported bounds are $\max(0,C-H)$ and $\min(1,C+H)$. They display
sampling uncertainty for this small case register; they do not make it
representative.

The one-sided permutation comparison preserves each method's flag count
per ticker, places flags among that ticker's observed dates without
replacement, and repeats this 1,000 times with seed 20261006. If $h_b$
is the simulated number of event hits and $h_{\rm obs}$ is observed,
\[
 p_{\rm perm}=\frac{1+\sum_{b=1}^{1000}{\bf1}(h_b\geq h_{\rm obs})}{1001}.
\]
The reference development value is 0.687 and held-out value is 1.000.
These values do not show a reliable advantage over random placement
under the current labels.

\section{Registered methods and baselines}
The robust variants use the previous 60 valid observations only. For
$M_t=\operatorname{median}(q_{<t})$ and
$\operatorname{MAD}_t=\operatorname{median}(|q-M_t|)$ over that window,
\[
 z^{\rm robust}_t=\frac{q_t-M_t}{\max(1.4826\,\operatorname{MAD}_t,f_q)}.
\]
The implementation applies cutoffs 3 and 4 to high log-volume or
absolute return scores. The floors are 0.10 for log-volume and 0.01
for returns. Market-adjusted return is the stock return minus exact-date
PX1 return, $r^{\rm excess}_{i,t}=r_{i,t}-r^{\rm PX1}_t$; its
prior-only mean/std score has a 0.01 denominator floor and absolute
cutoff 3.

The other registered candidates filter reference flags for positive
current volume, at least 20 active days and prior median volume at
least 400 over 60 observations; require another flag for the same symbol
within ten calendar days; or recompute the decay rule at month ends using
only information then available. The decay comparison uses preceding and
recent six-month windows, a prior mean volume of at least 1,000, a
recent/prior ratio below 0.20 or at least 90 days since positive volume.
The market-wide tag is context, not a seventh detector: it marked zero
of 1,399 company-share reference flags under its threshold.

Baselines are volume-only, absolute-return-only, market-adjusted
return-only, a simple volume spike, top two flags per month, and
same-count random placement. The latest-snapshot decay file was not
treated as a historical prediction. No candidate threshold was retuned.

\section{Development and held-out results}
\begin{table}[htbp]\centering\small
\caption{Selected primary-window comparisons; hits include 95\% Wilson intervals}
\begin{tabular}{@{}lccrr@{}}\toprule
Method & Development & Held out & Dev flags/week & Held flags/week\\\midrule
Reference v1 & 1/4 (0.05--0.70) & 0/1 (0.00--0.79) & 4.08 & 7.12\\
Return only & 1/4 (0.05--0.70) & 0/1 (0.00--0.79) & 3.54 & 6.30\\
Simple volume spike & 4/4 (0.51--1.00) & 1/1 (0.21--1.00) & 16.87 & 24.38\\
Robust z3 & 2/4 (0.15--0.85) & 0/1 (0.00--0.79) & 10.83 & 21.94\\
\bottomrule\end{tabular}
\end{table}
The simple volume-spike baseline hits all five eligible positives, but
its 16.87 and 24.38 flags per week are about four times the reference
workload of 4.08 and 7.12, and it produces 27 and 10 negative-control
flags versus the reference's 4 and 9; hit count alone is no advantage.
The complete method table remains in the generated Phase 2 comparison
output. The development-selected \code{robust_z3} hit zero of the one
held-out positive. Adoption required at least ten held-out positives,
a positive paired-bootstrap lower bound against both the reference and
the preselected best nonrandom baseline, and no workload or control
increase. The sample and comparative conditions fail. The decision is
to retain \code{reference_v1}. Across all five eligible positives,
the reference hits one in the primary window (Wilson 0.04--0.62) and two
in the secondary. None of these case-study rates estimates accuracy.

\section{Quality checks and verification}
The derived quality pass found zero duplicate symbol/date rows and zero
invalid dates, but 59 invalid numeric or OHLC rows and 349 gaps longer
than 30 calendar days. It marked 140 new-scrape rows for review:
134 with zero volume, five large jumps, and one OHLC inconsistency.
The corporate-action suspect list contains 115 rows and tags 86
reference flags. These are review prompts; no protected source CSV
was edited. The 100-headline human-label sample remains unfilled, so
no news-classifier accuracy or agreement is claimed.

The final offline suite passes 117 tests. The React production build
and Prettier format check pass. Playwright loaded the Terminal at 1280
and 390 pixels with no console errors and checked SPA refresh. The
committed frontend build is about 1.21 MB, allowing Flask to serve it
without Node. Verification establishes internal reproducibility, not
independent evidence of detector accuracy.
